% !TeX root = ../main.tex

\xchapter{结论与展望}{Conclusions and Future Work}

\xsection{结论}{Conclusions}

在视频数据指数级增长的背景下，视频摘要技术通过生成包含关键帧或片段的简洁摘要，为大规模视频的高效索引、浏览和检索提供了重要的技术支撑。然而，有监督方法受限于稀缺且主观性较强的标注数据，泛化能力不足；无监督方法主要基于底层视觉特征进行代表性筛选，在高层语义建模方面仍存在不足。因此，探索能够利用大语言模型通用语义知识、无需额外参数更新的免训练视频摘要方法具有显著的研究意义。其中，如何构建从视觉内容到文本语义空间的有效映射、如何设计多维度的重要性评估机制是关键所在。针对这些问题，本文提出了基于层级语义建模与全局主题约束的免训练视频摘要方法和基于多智能体协同推理与多维专家评分的视频摘要方法。本文的具体工作总结如下：

1）本文提出了基于层级语义建模与全局主题约束的免训练视频摘要方法。该方法包含视频场景划分、多级语义描述生成、全局主题约束的场景语义评分以及帧级语义分数增强四个核心模块。首先采用 KTS 算法将连续视频流组织为语义相近的场景单元；随后通过视觉语言模型生成帧级描述，并利用大语言模型沿帧级、场景级和视频级三个层次逐步聚合，将视觉内容映射到统一的文本语义空间；在此基础上，联合场景描述与视频级全局摘要评估各场景对整体主题的贡献程度；最后通过帧与场景、帧与全局描述之间的语义相关性计算，对场景级分数进行帧级增强，构建最终的帧级重要性序列。在 SumMe 和 TVSum 数据集上的实验表明，该方法在免训练方法中取得了最优的 $F_1$ 分数，在秩相关性评测中帧级排序与人工标注之间的一致性优于全部对比方法。消融实验进一步验证了帧级语义增强模块对整体性能的关键贡献。

2）本文提出了基于多智能体协同推理与多维专家评分的视频摘要方法。该方法构建了由 Planner Agent 与四类 Expert Agent 组成的多智能体协同推理框架，包含视频结构化处理与段级描述生成、全局规划与多专家协同评分、记忆增强与帧级摘要生成三个核心阶段。首先通过均匀分段与片段内多帧采样，利用视觉语言模型生成段级语义描述；随后由 Planner Agent 推断视频主题、生成全局摘要并动态分配专家权重，Story、Visual、Emotion 与 Information 四类 Expert Agent 分别从各自维度对片段进行独立评分，Planner Agent 综合各专家证据给出最终段级分数；最后通过轻量级历史记忆机制增强时序一致性，并将段级分数映射回帧级序列，经背包优化生成满足长度约束的最终摘要。实验结果表明，该方法在 SumMe 上的 $F_1$ 分数超越了全部监督方法基线，在秩相关性评测中两项排序指标均为全部对比方法中最优。智能体消融实验表明，Story Expert 与 Planner Agent 对整体性能贡献最为显著，各模块的边际贡献在不同数据集上呈现互补特征，印证了多专家协同评分的设计合理性。

\xsection{展望}{Future Work}

本文针对免训练视频摘要任务，分别提出了基于层级语义建模与全局主题约束的视频摘要方法和基于多智能体协同推理与多维专家评分的视频摘要方法。尽管这两种方法在公开数据集上取得了较好的性能，但在处理更长、内容更复杂的开放域视频时，仍在分段粒度适应性、长程语义建模能力以及面向用户需求的灵活性等方面存在不足，需要在未来的研究工作中加以改进：

1）自适应视频分段：本文两种方法分别采用 KTS 自适应划分与固定段数均匀切分，分段边界尚未充分利用视频内部的语义变化信息。未来可探索将内容感知的场景划分机制与多智能体框架相结合，使推理单元更贴近真实事件边界，提升段级描述与评分的一致性。

2）长程上下文建模：当前记忆机制通过拼接前序片段描述提供历史语境，对长程依赖与复杂叙事链的建模能力有限。未来可研究动态检索式的记忆组织机制，使模型能够针对当前片段选取更具判别价值的历史上下文，增强长视频场景下的全局一致性。

3）个性化视频摘要：本文方法面向通用摘要任务，尚未考虑用户的个性化偏好与查询需求。未来可通过动态调整专家权重或引入用户意图感知模块，将个性化信息融入多智能体评分框架，构建面向用户需求的视频摘要方法。